% ============================================================================
%  SECCIÓN 3.26: EVIDENCIAS DEL AVANCE
% ============================================================================

\section{Evidencias del avance}

A continuaci\'on se presentan las evidencias del avance correspondiente al
PDF~3. Las capturas de pantalla de cada ya fueron mostradas en la secci\'on de
construcci\'on de pantallas (secci\'on~3.17), por lo que aqu\'i se referencian
por su etiqueta y se complementan con evidencias adicionales del repositorio
y la ejecuci\'on en dispositivo.

% --- EVIDENCIA 1 ---
\subsubsection*{Evidencia 1: Pantallas implementadas}

Las cinco pantallas principales fueron implementadas y validadas en el
dispositivo f\'isico. Las capturas correspondientes se presentan en las
figuras~\ref{fig:pantalla-splash} a~\ref{fig:pantalla-perfil} de la
secci\'on~3.17.

\begin{itemize}
  \item Figura~\ref{fig:pantalla-splash}: Pantalla de bienvenida (Splash).
  \item Figura~\ref{fig:pantalla-login}: Pantalla de inicio de sesi\'on.
  \item Figura~\ref{fig:pantalla-registro}: Pantalla de registro.
  \item Figura~\ref{fig:pantalla-tabs}: Pantalla principal con barra de
    navegaci\'on inferior.
  \item Figura~\ref{fig:pantalla-scenario-detail}: Pantalla de detalle del
    escenario.
  \item Figura~\ref{fig:pantalla-perfil}: Pantalla de perfil de usuario.
\end{itemize}

% --- EVIDENCIA 2 ---
\subsubsection*{Evidencia 2: Navegaci\'on entre pantallas}

Se implement\'o la navegaci\'on completa entre pantallas utilizando
expo-router: Splash $\rightarrow$ Login $\rightarrow$ Registro, Login
$\rightarrow$ Tabs, y cierre de sesi\'on desde el perfil que regresa al
splash. El diagrama de navegaci\'on implementado se presenta en la
figura~\ref{fig:navegacion-implementada} de la secci\'on~3.21.

% --- EVIDENCIA 3 ---
\subsubsection*{Evidencia 3: Aplicaci\'on ejecutada en dispositivo f\'isico}

La aplicaci\'on se ejecut\'o y valid\'o en un tel\'efono f\'isico Android
mediante Expo Go, verificando el rendimiento, la interacci\'on y el
comportamiento de la interfaz en condiciones reales de uso.

\begin{figure}[H]
  \centering
  \caption{Evidencia: Aplicaci\'on ejecutada en dispositivo f\'isico}
  \label{fig:evidencia-dispositivo}
  \contenidoConFuente{%
    \includegraphics[width=0.48\linewidth]{evidencia-dispositivo-fisico.jpg}%
  }{Elaboraci\'on propia (fotograf\'ia del dispositivo f\'isico).}
\end{figure}

% --- EVIDENCIA 4 ---
\subsubsection*{Evidencia 4: Repositorio GitHub actualizado}

El repositorio del proyecto fue actualizado con los avances del PDF~3,
incluyendo las pantallas implementadas, la configuraci\'on del tema, los
componentes reutilizables y la navegaci\'on.

\begin{figure}[H]
  \centering
  \caption{Evidencia: Repositorio GitHub actualizado con avances del PDF 3}
  \label{fig:evidencia-github-pdf3}
  \contenidoConFuente{%
    \includegraphics[width=0.62\linewidth]{evidencia-repositorio-pdf3.png}%
  }{Elaboraci\'on propia (captura del repositorio en GitHub).}
\end{figure}

% --- EVIDENCIA 5 ---
\subsubsection*{Evidencia 5: Commits realizados}

Los commits realizados durante el avance documentan los principales hitos del
desarrollo:

\begin{lstlisting}[style=terminal]
feat(t007): implementar pantallas login y register con validacion
feat(nav): configurar flujo de navegacion y proteccion de rutas (T-010)
feat(ui): reemplazar emojis por Ionicons y agregar pantalla perfil con logout
feat: integrar splash screen con logo y boton de Comenzar
refactor(t012): extraer AuthGuard a componente reutilizable
\end{lstlisting}
